% ==STYLE-BEGIN==
\documentclass[UTF8,11pt]{ctexart}
\usepackage{amsmath,amssymb,amsthm}
\usepackage[a4paper,margin=2.4cm,top=2.3cm,bottom=2.3cm]{geometry}
\usepackage{fancyhdr}
\usepackage{enumitem}
\usepackage{hyperref}
\hypersetup{hidelinks}

% ---- 版式 ----
\linespread{1.15}
\setlength{\parskip}{0.35em}
\setlength{\parindent}{1.6em}
\setlist{nosep,leftmargin=2em}

% ---- 页眉 ----
\pagestyle{fancy}
\fancyhf{}
\fancyhead[L]{\small\coursename}
\fancyhead[R]{\small\thepage}
\renewcommand{\headrulewidth}{0.4pt}

% ---- 定理环境：中文用正体，标题加粗，按小节共用一条流水号 ----
\newtheoremstyle{cnthm}{0.7em}{0.7em}{\normalfont}{}{\bfseries}{}{0.5em}{}
\theoremstyle{cnthm}
\newtheorem{thm}{定理}[section]
\newtheorem{lem}[thm]{引理}
\newtheorem{prop}[thm]{命题}
\newtheorem{cor}[thm]{推论}
\newtheorem{defn}[thm]{定义}
\newtheorem{rem}[thm]{注}
\newtheorem{exa}[thm]{例}
\newtheorem{conj}[thm]{猜想}
\renewcommand{\proofname}{\bfseries 证明}
\providecommand{\coursename}{}
% ==STYLE-END==
\renewcommand{\coursename}{解析几何\,·\,0917}
\title{解析几何 0917}
\date{}
\begin{document}
\maketitle

\section{\ Recall}

\begin{itemize}
\item 概念：投影，线性映射，内积，水平分量，垂直分量；
\item 简化表述：克氏符号与点乘的坐标表示。
\end{itemize}

\section{\ 向量的外积（向量积，叉乘）}

\subsubsection*{右手系}

$\vec a,\vec b,\vec c$ 不共面（线性无关）。右手四指从 $\vec a$
转到 $\vec b$，如果大拇指所指方向与 $\vec c$ 一致，则称
$(\vec a,\vec b,\vec c)$ 是右手系，否则是左手系。

若 $(\vec a,\vec b,\vec c)$ 是右手系，则 $\vec a,\vec b,\vec c$
中任意两个向量变换位置，则得到新的坐标系是左手系。

\subsubsection*{叉乘}

$\vec a,\vec b$ 是两个非零向量，则 $\vec a,\vec b$ 的外积
$\vec a\times\vec b$ 定义为

\begin{itemize}
\item 长度为 $|\vec a||\vec b|\sin\langle\vec a,\vec b\rangle$；
\item 方向为与 $\vec a,\vec b$ 垂直，且
$(\vec a,\vec b,\vec a\times\vec b)$ 构成右手系。
\end{itemize}

特别地，如果 $\vec a$ or $\vec b=\vec 0$，或者
$\vec a,\vec b$ 共线，则 $\vec a\times\vec b=\vec 0$。

\subsubsection*{外积的几何意义}

\begin{enumerate}
\item $|\vec a\times\vec b|
=|\vec a||\vec b|\sin\langle\vec a,\vec b\rangle$，
它表示 $\vec a,\vec b$ 围成的平行四边形的面积；
\item $\vec e\times\vec a=\vec e\times\vec a_h$，
即 $\vec e\times\vec a$ 是将 $\vec a$ 的外支量绕 $\vec e$
方向右旋 $\dfrac\pi2$。

$\vec a$ 绕 $\vec e$ 方向右旋 $\theta$ 角，意味着将其外支量绕
$\vec e$ 右旋 $\theta$ 角。
\end{enumerate}

\subsubsection*{外积的性质}

\begin{enumerate}
\item \textbf{反交换律}\quad
$\vec a\times\vec b=-\vec b\times\vec a$。

$(\vec a,\vec b,\vec a\times\vec b)$ 与
$(\vec a,\vec b,\vec b\times\vec a)$ 一定一个为左手系，
一个为右手系。
\item \textbf{结合律}\quad
$(\lambda\vec a)\times\vec b=\lambda(\vec a\times\vec b)$。
\item \textbf{分配律}\quad
$\vec e\times(\vec a+\vec b)
=\vec e\times\vec a+\vec e\times\vec b$。
\end{enumerate}

\end{document}
